% ============================================================
% DADOS: ANEXO J - NOTIFICAÇÃO DE ENVOLVIDO QUE PASSA À CONDIÇÃO DE SINDICADO
% Arquivo: dados/dados-anexo-j-notificacao-novo-sindicado.tex
% ============================================================

% ------------------------------------------------------------
% Data própria do documento
% Deixe em branco para usar a data/local do termo de abertura.
% Formato: \dataDocumento{dia}{mês}{ano}{cidade/UF}
% ------------------------------------------------------------
\dataDocumento{6}{maio}{2026}{CIDADE/UF}
\chaveCronologica{20260506}

\novaNotificacaoNovoSindicado

\notificacaoNovoSindicadoNumero{001}
\notificacaoNovoSindicadoOrigem{Sindicante}

% Se ficar vazio, o modelo usa \localData.
% \notificacaoNovoSindicadoLocalData{CIDADE/UF, 31 de maio de 2026}

\notificacaoNovoSindicadoDestinatario{\getSindicadoPosto\ \getSindicadoNome\ -- \getSindicadoSU/\getSindicadoOM}
\notificacaoNovoSindicadoAnexos{cópia da \getPortaria; cópia dos documentos que deram origem à instauração}
\notificacaoNovoSindicadoLocalVista{Sala do Sindicante, no quartel do \getOm}
\notificacaoCanal{direta}
\notificacaoAntecedenciaDiasUteis{3}
\notificacaoFatosEspecificos{elementos apurados que indicam a possível participação do destinatário nos fatos delimitados pela portaria instauradora}
\notificacaoFormaCiencia{recibo pessoal datado e assinado ao final deste documento}

% Campos opcionais, se houver reinquirição já marcada.
\notificacaoNovoSindicadoDataReinquiricao{00 de mês de 2026}
\notificacaoNovoSindicadoHoraReinquiricao{00h00}
\notificacaoNovoSindicadoLocalReinquiricao{Sala do Sindicante, no quartel do \getOm}
